% !TEX program = pdflatex
% DS Physique-Chimie — 2nde Bac Pro — Chapitre 10
% Température et capteurs thermiques
% Durée : 30 minutes — Barème : /20
% Notions évaluées :
%   - Conversion Celsius <-> Kelvin (dans les deux sens)
%   - Exploitation d'une courbe d'étalonnage de capteur (CTN, thermocouple)

\documentclass[a4paper,11pt]{article}
\usepackage[utf8]{inputenc}
\usepackage[T1]{fontenc}
\usepackage[french]{babel}
\usepackage{geometry}
\geometry{left=1.8cm,right=1.8cm,top=1.6cm,bottom=1.8cm}
\usepackage{amsmath,amssymb}
\usepackage{enumitem}
\usepackage{array}
\usepackage{booktabs}
\usepackage{xcolor}
\usepackage{tcolorbox}
\usepackage{fancyhdr}
\usepackage{lastpage}
\usepackage{tikz}
\usepackage{pgfplots}
\pgfplotsset{compat=1.17}

\tcbuselibrary{skins,breakable}

% Couleurs
\definecolor{violet}{HTML}{6f42c1}
\definecolor{violetclair}{HTML}{f5f0ff}
\definecolor{orange}{HTML}{d97706}
\definecolor{bleu}{HTML}{0056b3}
\definecolor{vert}{HTML}{2da44e}
\definecolor{rouge}{HTML}{c53030}
\definecolor{gris}{HTML}{718096}

% En-tête / pied
\pagestyle{fancy}
\fancyhf{}
\lhead{\small DS Physique-Chimie — 2nde Bac Pro}
\rhead{\small Ch.\,10 — Température et capteurs}
\cfoot{\thepage\,/\,\pageref{LastPage}}
\renewcommand{\headrulewidth}{0.5pt}

% Boîtes
\newtcolorbox{enonce}{
  colback=violetclair,
  colframe=violet,
  leftrule=4pt,
  rightrule=0pt,toprule=0pt,bottomrule=0pt,
  boxsep=4pt,left=8pt,right=6pt,
  breakable
}
\newtcolorbox{rappel}{
  colback=yellow!6,
  colframe=orange,
  title={\textbf{Rappels utiles}},
  fonttitle=\bfseries,
  breakable
}
\newtcolorbox{consigne}{
  colback=blue!3,
  colframe=bleu,
  title={\textbf{Consignes}},
  fonttitle=\bfseries
}

% Marge pour barème
\newcommand{\pt}[1]{\hfill\textcolor{gris}{\small\textit{/\,#1 pt}}}

\begin{document}

%===============================================================
% EN-TÊTE
%===============================================================
\begin{center}
{\LARGE\bfseries\color{violet} Devoir Surveillé — Physique-Chimie}\\[4pt]
{\large\bfseries Chapitre 10 — Température et capteurs thermiques}\\[4pt]
{\normalsize\itshape Atelier de menuiserie — Filière bois}\\[6pt]
{\normalsize Seconde Baccalauréat Professionnel — Durée : \textbf{30 minutes} — Barème : \textbf{/20}}
\end{center}

\vspace{2pt}
\noindent
\begin{tabular}{@{}p{0.55\textwidth}p{0.4\textwidth}@{}}
\textbf{Nom :} \dotfill & \textbf{Prénom :} \dotfill \\[6pt]
\textbf{Classe :} \dotfill & \textbf{Date :} \dotfill \\
\end{tabular}

\vspace{4pt}
\begin{consigne}
\begin{itemize}[leftmargin=*,itemsep=1pt,topsep=2pt]
  \item Calculatrice autorisée. Téléphone interdit.
  \item Soigner la présentation, poser les calculs, indiquer les unités.
  \item Les réponses non justifiées ne seront pas prises en compte.
\end{itemize}
\end{consigne}

\begin{rappel}
\begin{itemize}[leftmargin=*,itemsep=1pt,topsep=2pt]
  \item \textbf{Conversion des températures :} $T\,(\text{K}) = \theta\,(^\circ\text{C}) + 273$
        \quad et \quad $\theta\,(^\circ\text{C}) = T\,(\text{K}) - 273$
  \item \textbf{Zéro absolu :} $0$\,K $= -273\ ^\circ$C (température la plus basse possible).
  \item \textbf{Courbe d'étalonnage :} relie la grandeur électrique mesurée (résistance $R$ en $\Omega$, tension $U$ en mV) à la température réelle.
  \item \textbf{Interpolation linéaire} entre deux points $(T_1,y_1)$ et $(T_2,y_2)$ :
        \[ y = y_1 + \frac{T - T_1}{T_2 - T_1}\,(y_2 - y_1) \]
\end{itemize}
\end{rappel}

%===============================================================
% EXERCICE 1 — CONVERSIONS
%===============================================================
\vspace{4pt}
\noindent{\large\bfseries\color{violet}Exercice 1 — Du Celsius au Kelvin et inversement \pt{6}}

\begin{enonce}
Dans un atelier de menuiserie, un technicien intervient sur plusieurs équipements qui nécessitent un contrôle précis de la température : séchoir à bois, local de stockage des colles, chaudière à copeaux\ldots{} Il doit remplir une fiche de maintenance dans laquelle les températures doivent être exprimées en Kelvin, alors que ses instruments affichent des degrés Celsius (et inversement).
\end{enonce}

\begin{enumerate}[label=\textbf{\arabic*.},leftmargin=*,itemsep=4pt]

\item Le séchoir à bois est réglé sur la consigne $\theta = 65\ ^\circ$C pour le séchage de planches de chêne. Exprimer cette température en Kelvin. \pt{1}

\vspace{1.2cm}

\item La température ambiante de l'atelier de menuiserie est $\theta = 18\ ^\circ$C. Convertir cette température en Kelvin. \pt{1}

\vspace{1.2cm}

\item Une sonde Pt100 placée dans le conduit de sortie d'une \textbf{chaudière à copeaux de bois} indique $T = 313$\,K. Exprimer cette température en degrés Celsius. \pt{1}

\vspace{1.2cm}

\item Le local de stockage des \textbf{colles vinyliques} pour panneaux de menuiserie est maintenu à $T = 268$\,K afin de préserver leur qualité. Convertir cette température en degrés Celsius. \pt{1}

\vspace{1.2cm}

\item Compléter le tableau de correspondance ci-dessous. \pt{1,5}

\vspace{2pt}
\begin{center}
\renewcommand{\arraystretch}{1.4}
\begin{tabular}{|>{\columncolor{violetclair}}c|c|c|c|c|c|}
\hline
\rowcolor{violetclair}
$\theta\ (^\circ$C) & $-273$ & $0$ & \rule{1.3cm}{0.3pt} & $100$ & \rule{1.3cm}{0.3pt} \\
\hline
$T\ $(K) & \rule{1.3cm}{0.3pt} & \rule{1.3cm}{0.3pt} & $293$ & \rule{1.3cm}{0.3pt} & $573$ \\
\hline
\end{tabular}
\end{center}

\vspace{4pt}

\item \textbf{Justifier} pourquoi il est impossible de mesurer une température de $-10$\,K dans l'atelier de menuiserie. \pt{0,5}

\vspace{1.2cm}
\end{enumerate}

\newpage

%===============================================================
% EXERCICE 2 — EXPLOITATION COURBE CTN
%===============================================================
\noindent{\large\bfseries\color{violet}Exercice 2 — Courbe d'étalonnage de la CTN d'un séchoir à bois \pt{9}}

\begin{enonce}
Dans un \textbf{séchoir à bois} d'un atelier de menuiserie, la température est mesurée par une \textbf{thermistance CTN} afin de réguler le chauffage et obtenir un séchage homogène des planches. Le tableau et la courbe d'étalonnage ci-dessous donnent la résistance $R$ de la CTN en fonction de la température $\theta$ à l'intérieur du séchoir.
\end{enonce}

\begin{center}
\renewcommand{\arraystretch}{1.3}
\begin{tabular}{|>{\columncolor{violetclair}}c|c|c|c|c|c|}
\hline
\rowcolor{violetclair}
$\theta\ (^\circ$C) & $20$ & $40$ & $60$ & $80$ & $100$ \\
\hline
$R\ (\Omega)$ & $10\,000$ & $5\,000$ & $2\,000$ & $1\,000$ & $500$ \\
\hline
\end{tabular}
\end{center}

\vspace{4pt}
\begin{center}
\begin{tikzpicture}
\begin{axis}[
  width=17cm, height=12cm,
  xlabel={\textbf{Température $\theta$ (en $^\circ$C)}},
  ylabel={\textbf{Résistance $R$ (en $\Omega$)}},
  xlabel style={font=\small},
  ylabel style={font=\small},
  xmin=10, xmax=110,
  ymin=0, ymax=11000,
  xtick={10,20,30,40,50,60,70,80,90,100,110},
  minor x tick num=4,
  ytick={0,500,1000,1500,2000,2500,3000,3500,4000,4500,5000,5500,6000,6500,7000,7500,8000,8500,9000,9500,10000,10500,11000},
  yticklabel style={font=\footnotesize},
  xticklabel style={font=\footnotesize},
  scaled y ticks=false,
  y tick label style={/pgf/number format/fixed,/pgf/number format/use comma},
  minor y tick num=4,
  grid=both,
  major grid style={line width=.5pt,draw=gray!55},
  minor grid style={line width=.25pt,draw=gray!25},
  title={\textbf{Courbe d'étalonnage de la CTN du séchoir à bois}},
  title style={font=\small},
  every axis plot/.append style={very thick},
]
\addplot[violet,smooth,mark=*,mark size=2.6pt,mark options={fill=violet}]
  coordinates {(20,10000) (40,5000) (60,2000) (80,1000) (100,500)};
\end{axis}
\end{tikzpicture}
\end{center}

\vspace{2pt}
\begin{center}\small\itshape Utiliser la courbe ci-dessus pour répondre aux questions suivantes (lecture graphique).\end{center}

\newpage
\noindent{\bfseries\color{violet}Exercice 2 (suite) — Questions}

\vspace{4pt}
\begin{enumerate}[label=\textbf{\arabic*.},leftmargin=*,itemsep=4pt]

\item Comment évolue la résistance $R$ de la CTN lorsque la température $\theta$ augmente ? En déduire la signification du sigle \textbf{CTN}. \pt{1,5}

\vspace{1.4cm}

\item L'ohmmètre branché sur la CTN affiche $R = 5\,000\ \Omega$. En utilisant la courbe d'étalonnage, déterminer la température $\theta$ à l'intérieur du séchoir à bois. \pt{1,5}

\vspace{1.4cm}

\item La consigne de séchage pour des planches de hêtre est fixée à $\theta = 80\ ^\circ$C. Quelle est la valeur de la résistance $R$ que doit détecter le régulateur du séchoir ? \pt{1,5}

\vspace{1.4cm}

\item Convertir cette température de consigne ($80\ ^\circ$C) en Kelvin. \pt{1}

\vspace{1.2cm}

\item Pendant un cycle de séchage, on mesure une résistance $R = 3\,500\ \Omega$. Cette valeur ne figure pas dans le tableau.\par
\textbf{a.} Entre quelles deux valeurs de température se situe $\theta$ ? \pt{1}

\vspace{1.2cm}

\textbf{b.} En utilisant l'\textbf{interpolation linéaire} entre $\theta_1 = 40\ ^\circ$C ($R_1 = 5\,000\ \Omega$) et $\theta_2 = 60\ ^\circ$C ($R_2 = 2\,000\ \Omega$), calculer la valeur de la température mesurée. \pt{2,5}

\vspace{2.2cm}
\end{enumerate}

\newpage
%===============================================================
% EXERCICE 3 — THERMOCOUPLE
%===============================================================
\noindent{\large\bfseries\color{violet}Exercice 3 — Thermocouple dans une chaudière à copeaux \pt{5}}

\begin{enonce}
L'atelier de menuiserie est chauffé par une \textbf{chaudière à copeaux de bois} qui brûle les déchets issus des machines (scie à panneaux, raboteuse, dégauchisseuse). Dans la chambre de combustion, la température est mesurée par un \textbf{thermocouple de type K}. La tension $U$ délivrée (en mV) est proportionnelle à la température $\theta$ (en $^\circ$C), selon la relation :
\[ U\ (\text{mV}) = 0{,}04 \times \theta\ (^\circ\text{C}) \]
\end{enonce}

\begin{enumerate}[label=\textbf{\arabic*.},leftmargin=*,itemsep=4pt]

\item Au démarrage de la chaudière, la température dans la chambre de combustion est $\theta = 250\ ^\circ$C. Calculer la tension $U$ délivrée par le thermocouple. \pt{1,5}

\vspace{1.6cm}

\item En régime de fonctionnement normal, le voltmètre branché sur le thermocouple affiche $U = 16$\,mV. En déduire la valeur de la température $\theta$ dans la chambre de combustion, en degrés Celsius. \pt{1,5}

\vspace{1.6cm}

\item Exprimer cette température en Kelvin. \pt{1}

\vspace{1.2cm}

\item En pleine combustion, la chambre peut atteindre $900\ ^\circ$C. Parmi les capteurs suivants : \emph{CTN}, \emph{Pt100}, \emph{thermocouple type K}, lequel est le plus adapté pour mesurer une telle température ? Justifier brièvement. \pt{1}

\vspace{1.6cm}
\end{enumerate}

\vspace{4pt}
\hrulefill\\[2pt]
\hfill\textit{\small Fin du sujet — Vérifier ses réponses avant de rendre la copie.}

\end{document}
